\documentclass[11pt]{article}
\usepackage[margin=2.5cm]{geometry}
\usepackage{amsmath,amssymb,amsthm}
\usepackage{booktabs}
\usepackage{graphicx}
\usepackage{hyperref}

\newtheorem{theorem}{Theorem}
\newtheorem{proposition}[theorem]{Proposition}
\newtheorem{observation}[theorem]{Observation}
\theoremstyle{definition}
\newtheorem{definition}[theorem]{Definition}
\theoremstyle{remark}
\newtheorem{remark}[theorem]{Remark}

\title{\textbf{A Martin-L\"of Test That\\
Pseudorandom Generators Fail}}
\author{Robin Langer}
\date{August 2026}

\begin{document}
\maketitle

\begin{abstract}
Every computable sequence fails some Martin-L\"of test,
but the standard proof is circular:
it uses the sequence itself to build the test.
We construct a concrete, non-circular Martin-L\"of test---based
on the Chebotarev density theorem---that five standard
pseudorandom number generators all fail.
The test asks whether a bitstream indexed by primes
has the correct conditional hit rates modulo~$12$.
Any sequence drawn from the Chebotarev measure passes.
Any PRNG fails, because its bits are generated
independently of the primes.
\end{abstract}

%% ================================================================
\section{The problem with the standard proof}
%% ================================================================

A pseudorandom number generator produces a computable sequence.
No computable sequence is Martin-L\"of random.
The proof: given a computable sequence~$x$,
define $U_m = \{y : y_i = x_i \text{ for } i \leq f(m)\}$
where $f(m)$ grows slowly enough that $\mu(U_m) \leq 2^{-m}$.
Then $x \in \bigcap U_m$, so $x$ fails the test.

This is correct but useless.
The test requires you to \emph{compute $x$} before you can test it.
For a PRNG with a short seed, that is feasible.
But for a number-theoretic sequence like the Artin membership
sequence, the test requires factoring $q - 1$ for every prime~$q$---the
test is harder than the sequence it claims to debunk.

More importantly, the test is not \emph{interpretable}.
It tells you that the sequence is not random
but says nothing about \emph{how} it fails.
We construct a test that is concrete,
interpretable, and motivated by number theory.

%% ================================================================
\section{The Chebotarev measure}
%% ================================================================

Fix an integer~$r$.
A prime~$q$ belongs to $S(r)$ if $r$ generates
the multiplicative group $(\mathbb{Z}/q\mathbb{Z})^*$.
Artin conjectured that $S(r)$ has density
$C_{\mathrm{Artin}} \approx 0.3740$ among the primes.

But the density is not uniform across residue classes.
Whether $r$ is a primitive root mod~$q$ depends on
the factorisation of $q - 1$,
which is constrained by $q \bmod m$ for small~$m$.
For $r = 3$ and $m = 12$:

\begin{itemize}
\item $q \equiv 1 \pmod{12}$: then $3 \mid q - 1$,
  so the cubes are a proper subgroup.
  $3$ cannot generate. Hit rate $= 0$.
\item $q \equiv 11 \pmod{12}$: again $3 \mid q - 1$. Hit rate $= 0$.
\item $q \equiv 5 \pmod{12}$: neither $3$ nor $2$
  divides $q - 1$ obstructively. Hit rate $\approx 0.90$.
\item $q \equiv 7 \pmod{12}$: $3 \nmid q - 1$ but $2 \| q - 1$.
  Hit rate $\approx 0.60$.
\end{itemize}

These are not empirical observations.
The vanishing at $q \equiv 1$ and $q \equiv 11$ is a theorem.
The values $0.90$ and $0.60$ are predictions of the
Chebotarev density theorem applied to the splitting field
of $x^{q-1} - 1$ over $\mathbb{Q}$.

\begin{definition}
The \emph{Chebotarev measure} $\mu$ on $\{0,1\}^{\mathbb{N}}$,
indexed by the primes $p_1, p_2, \ldots$, is defined by
\[
  \mu(b_n = 1 \mid p_n \equiv a \pmod{12})
  = \begin{cases}
    0 & \text{if } a \in \{1, 11\}, \\
    c_a & \text{if } a \in \{5, 7\},
  \end{cases}
\]
where $c_5 \approx 0.897$ and $c_7 \approx 0.600$
are the Chebotarev-predicted densities,
with bits conditionally independent given the residue class.
\end{definition}

This is a computable measure:
the residue class $p_n \bmod 12$ is computable,
and the conditional probabilities are computable reals.
Martin-L\"of randomness with respect to $\mu$ is well-defined.

%% ================================================================
\section{The test}
%% ================================================================

\begin{definition}
The \emph{residue-class coherence test} is the
Martin-L\"of test $(U_m)_{m \geq 1}$ relative to $\mu$ defined by
\[
  U_m = \left\{ (b_n) :
    \sum_{a \in \{1,5,7,11\}}
    \frac{(\hat{P}_a - \mu_a)^2}{\mu_a(1-\mu_a) / N_a}
    > t_m \right\},
\]
where $\hat{P}_a$ is the empirical hit rate
among the first $N$ primes congruent to $a \pmod{12}$,
$\mu_a$ is the Chebotarev prediction,
$N_a = |\{n \leq N : p_n \equiv a \pmod{12}\}|$,
and $t_m$ is chosen so that $\mu(U_m) \leq 2^{-m}$.
\end{definition}

Under $\mu$, the statistic converges in distribution
to $\chi^2_3$ as $N \to \infty$,
so $t_m = O(m)$ suffices.
The sets $U_m$ are uniformly computably enumerable
(the chi-squared statistic is computable from the data),
and $\mu(U_m) \leq 2^{-m}$ by choice of $t_m$.
This is a valid Martin-L\"of test.

%% ================================================================
\section{Pseudorandom generators fail}
%% ================================================================

A PRNG generates bits independently of the prime index.
Its bits satisfy
$P(b_n = 1 \mid p_n \equiv a \pmod{12}) = P(b_n = 1)$
for all~$a$.  In particular, $\hat{P}_1 \approx \hat{P}_5
\approx \hat{P}_7 \approx \hat{P}_{11} \approx 0.375$,
while $\mu_1 = \mu_{11} = 0$ and $\mu_5 \approx 0.90$.

As $N \to \infty$, the chi-squared statistic diverges.
The PRNG output lies in $U_m$ for every~$m$.
It fails the test.

We verified this for five standard generators
over $N = 50{,}000$ primes:

\begin{center}
\begin{tabular}{lccccc}
\toprule
Generator & $\hat{P}_1$ & $\hat{P}_5$ & $\hat{P}_7$
  & $\hat{P}_{11}$ & $\chi^2$ \\
\midrule
Mersenne Twister   & 0.369 & 0.370 & 0.380 & 0.382 & 4.6 \\
PCG64              & 0.378 & 0.372 & 0.379 & 0.376 & 0.7 \\
MT19937            & 0.375 & 0.372 & 0.376 & 0.367 & 1.6 \\
Philox             & 0.377 & 0.368 & 0.369 & 0.378 & 3.0 \\
SFC64              & 0.374 & 0.377 & 0.371 & 0.379 & 1.2 \\
\bottomrule
\end{tabular}
\end{center}

The chi-squared values above are computed against the
\emph{flat} null ($\hat{P}_a = \hat{P}$ for all~$a$).
Against the Chebotarev measure $\mu$,
the statistic is dominated by the $a = 1$ and $a = 11$
terms, where $\mu_a = 0$ but $\hat{P}_a \approx 0.375$,
and formally diverges.
The generators do not merely fail the test.
They fail it infinitely badly.

%% ================================================================
\section{The Artin sequence passes}
%% ================================================================

The Artin membership sequence
$b_n = \mathbf{1}_{p_n \in S(3)}$, computed over $50{,}000$ primes:

\begin{center}
\begin{tabular}{rccc}
\toprule
$p_n \bmod 12$ & $N_a$ & Hits & $\hat{P}_a$ \\
\midrule
$1$  & $12{,}447$ & $0$      & $0.000$ \\
$5$  & $12{,}538$ & $11{,}243$ & $0.897$ \\
$7$  & $12{,}510$ & $7{,}508$  & $0.600$ \\
$11$ & $12{,}503$ & $0$      & $0.000$ \\
\bottomrule
\end{tabular}
\end{center}

The empirical hit rates match the Chebotarev predictions exactly.
The chi-squared statistic against $\mu$ is small.
The Artin sequence passes the residue-class coherence test.

Note that this does not make the Artin sequence Martin-L\"of random.
It is computable, so it fails other tests.
But it passes \emph{this} test---the one that every PRNG fails.

%% ================================================================
\section{The Artin sequence fails standard tests}
%% ================================================================

At the same time, the Artin bitstream fails
five of nine standard PRNG tests:

\begin{center}
\begin{tabular}{lrl}
\toprule
Test & $p$-value & Verdict \\
\midrule
Frequency (monobit vs $0.5$) & $\approx 0$ & Fail \\
Frequency (vs Artin's constant) & $0.565$ & Pass \\
Runs & $\approx 0$ & Fail \\
Serial (bigrams) & $\approx 0$ & Fail \\
Autocorrelation (lags $1$--$50$) & $\approx 0$ & Fail \\
Block frequency ($M=100$) & $1.000$ & Pass \\
Longest run of ones & $0.018$ & Pass \\
CUSUM & $\approx 0$ & Fail \\
Spectral (DFT) & $0.296$ & Pass \\
\bottomrule
\end{tabular}
\end{center}

The most informative failure is the lag-1 autocorrelation
$r_1 = -0.062$: consecutive primes repel each other in $S(3)$.
\[
  P(p_{n+1} \in S(3) \mid p_n \in S(3)) = 0.337,
  \qquad
  P(p_{n+1} \in S(3) \mid p_n \notin S(3)) = 0.398.
\]
This repulsion varies with $r$ (Table~\ref{tab:repulsion}).
Most values repel, but $r = 15$ attracts
and $r = 14$ is neutral.

\begin{table}[h]
\centering
\begin{tabular}{rccccr}
\toprule
$r$ & $P$ & $P|\text{hit}$ & $P|\text{miss}$ & Drop & $r_1$ \\
\midrule
2  & 0.374 & 0.329 & 0.401 & $+12.0\%$ & $-0.072$ \\
3  & 0.375 & 0.337 & 0.398 & $+10.3\%$ & $-0.062$ \\
5  & 0.394 & 0.338 & 0.431 & $+14.3\%$ & $-0.093$ \\
6  & 0.374 & 0.351 & 0.387 & $+6.2\%$  & $-0.037$ \\
7  & 0.374 & 0.369 & 0.377 & $+1.3\%$  & $-0.008$ \\
10 & 0.376 & 0.362 & 0.384 & $+3.6\%$  & $-0.021$ \\
11 & 0.375 & 0.359 & 0.385 & $+4.4\%$  & $-0.026$ \\
12 & 0.376 & 0.333 & 0.401 & $+11.4\%$ & $-0.069$ \\
13 & 0.377 & 0.346 & 0.396 & $+8.3\%$  & $-0.050$ \\
14 & 0.374 & 0.374 & 0.373 & $-0.1\%$  & $+0.000$ \\
\textbf{15} & \textbf{0.377} & \textbf{0.420} & \textbf{0.351}
  & $\mathbf{-11.4\%}$ & $\mathbf{+0.069}$ \\
17 & 0.376 & 0.350 & 0.391 & $+6.8\%$  & $-0.041$ \\
18 & 0.373 & 0.329 & 0.399 & $+11.7\%$ & $-0.070$ \\
19 & 0.376 & 0.373 & 0.378 & $+0.7\%$  & $-0.004$ \\
20 & 0.394 & 0.335 & 0.432 & $+14.8\%$ & $-0.096$ \\
\bottomrule
\end{tabular}
\caption{Consecutive-prime repulsion in $S(r)$.
Positive drop means repulsion; negative means attraction.}
\label{tab:repulsion}
\end{table}

%% ================================================================
\section{Discussion}
%% ================================================================

The standard proof that computable sequences are not Martin-L\"of
random is an existence result:
for every computable sequence, \emph{some} test catches it.
The test is built from the sequence, cannot be described
without it, and conveys no insight about
\emph{why} the sequence is not random.

We have exhibited a specific test---residue-class coherence
against the Chebotarev measure---with three properties:
\begin{enumerate}
\item It is a valid Martin-L\"of test
  (computably enumerable, measure-decreasing).
\item It catches every PRNG
  (the bits are independent of the prime index).
\item It has a mathematical name and an interpretable failure mode
  (the generator does not know which residue class it is in).
\end{enumerate}

The test does not use the PRNG's algorithm or seed.
It works against \emph{any} generator that produces bits
independently of the primes,
which includes all standard PRNGs.

The Artin sequence is not Martin-L\"of random either---it
is computable, and other tests catch it.
But the tests that catch it (autocorrelation, runs)
are \emph{different} from the test that catches PRNGs.
The two kinds of non-randomness do not overlap.

\medskip

The situation is:
\begin{center}
\begin{tabular}{lcc}
\toprule
& Standard tests & Coherence test \\
\midrule
Truly random sequence & Pass & Pass \\
Artin sequence & Fail & Pass \\
PRNG & Pass & Fail \\
\bottomrule
\end{tabular}
\end{center}

\medskip

No single test captures ``randomness.''
What exists is a landscape of tests,
each sensitive to a different kind of structure.
Standard PRNG tests detect the absence of independence.
The coherence test detects the absence of arithmetic.
The Artin sequence has arithmetic but not independence.
A PRNG has independence but not arithmetic.
A truly random sequence---drawn from the Chebotarev measure---would
have both.

\begin{thebibliography}{9}
\bibitem{artin1927}
E.~Artin,
\emph{\"Uber die Zetafunktionen gewisser algebraischer Zahlk\"orper},
Math.\ Ann.\ \textbf{89} (1927), 147--156.

\bibitem{hooley1967}
C.~Hooley,
\emph{On Artin's conjecture},
J.~Reine Angew.\ Math.\ \textbf{225} (1967), 209--220.

\bibitem{heathbrown1986}
D.R.~Heath-Brown,
\emph{Artin's conjecture for primitive roots},
Quart.\ J.\ Math.\ Oxford \textbf{37} (1986), 27--38.

\bibitem{martinlof1966}
P.~Martin-L\"of,
\emph{The definition of random sequences},
Information and Control \textbf{9}(6) (1966), 602--619.

\bibitem{nist2010}
A.~Rukhin et al.,
\emph{A Statistical Test Suite for Random and Pseudorandom
Number Generators for Cryptographic Applications},
NIST SP~800-22 Rev.~1a (2010).

\bibitem{lecuyer2007}
P.~L'Ecuyer and R.~Simard,
\emph{TestU01: A C library for empirical testing of
random number generators},
ACM Trans.\ Math.\ Software \textbf{33}(4) (2007), Article~22.
\end{thebibliography}

\end{document}
